\documentclass{article}%
\usepackage[T1]{fontenc}%
\usepackage[utf8]{inputenc}%
\usepackage{lmodern}%
\usepackage{textcomp}%
\usepackage{lastpage}%
\usepackage{geometry}%
\geometry{margin=2cm}%
\usepackage{expex}%
\usepackage[T1]{fontenc}%
\usepackage[utf8]{inputenc}%
\usepackage[spanish]{babel}%
%
\title{Análisis Lingüístico: Gramatica Descriptiva Qanjobal}%
\author{Generado por el TFM de Elías Carrasco}%
%
\begin{document}%
\normalsize%
\maketitle%
\section{Page 97}%
\label{sec:Page97}%
\subsection{4.2 PRONOMBRES}%
\label{subsec:4.2PRONOMBRES}%

%
Los pronombres son de diversos tipos, y cada grupo está formado por distintos %
morfemas. Lo que tienen en común es que siempre llevan el morfema de persona %
perteneciente al Juego Absolutivo. De igual manera, cada tipo de pronombre corresponde %
a un nominal pero con funciones y ambientes específicos. %
\subsection{4.2.1 Clasificación de pronombres}%
\label{subsec:4.2.1Clasificacindepronombres}%

%
Los tipos de pronombres son: pronombres independientes, pronombres ergativos y %
pronombres absolutivos. %
\subsection{4.2.1.1 Pronombres independientes}%
\label{subsec:4.2.1.1Pronombresindependientes}%

%
Los pronombres independientes, están formados básicamente por los morfemas %
partículas enfática + Juego B. Estos pronombres tienen variaciones en su %
pronunciación pero, en general su estructura es la indicada. %
\subsection{4.2.1.2 Pronombres ergativos}%
\label{subsec:4.2.1.2Pronombresergativos}%

%
Los pronombres ergativos son los mismos marcadores de Juego A que ya se %
explicó anteriormente. Sus funciones son tres: poseedor de sustantivos, sujeto de %
verbos transitivos y sujeto de verbos intransitivos en ergatividad mixta. Hay %
pronombres ergativos preconsonánticos (pronombres usados con palabras que inician %
con consonante), y pronombres prevocálicos (morfemas de persona que se usan con %
palabras que inician con vocal). El paradigma completo de estos pronombres es el %
siguiente: 

%
\end{document}